\begin{frame}{Partial Scan Creates an X-State Barrier}
  \begin{itemize}
    \item Scan-MUX delay breaks timing $\to$ top \alert{10\%} FFs kept \alert{non-scan}
    \item At one frame, non-scan FF outputs are \alert{unknown (X)}
  \end{itemize}
  \vspace{0.5em}
  \centering
  % X-state blocks propagation: fault -> logic -> non-scan FF(X) --X--> output
\begin{tikzpicture}[
  >=Stealth, node distance=0.85cm,
  blk/.style={draw=PubBlue, fill=PubBlue!80, text=black, rounded corners=3pt,
              minimum width=1.7cm, minimum height=0.85cm, align=center, font=\scriptsize\bfseries},
  xff/.style={draw=WarnRed, fill=WarnRed!85, text=black, rounded corners=3pt,
              minimum width=2.0cm, minimum height=0.85cm, align=center, font=\scriptsize\bfseries},
  obs/.style={draw=GoodGreen, fill=GoodGreen!80, text=black, rounded corners=3pt,
              minimum width=1.6cm, minimum height=0.85cm, align=center, font=\scriptsize\bfseries},
  arr/.style={->, thick, gray!70},
]
  \node[blk] (f) {fault};
  \node[blk, right=of f] (l) {logic};
  \node[xff, right=of l] (ff) {non-scan FF\\= X};
  \node[obs, right=1.5cm of ff] (o) {output};
  \draw[arr] (f) -- (l);
  \draw[arr] (l) -- (ff);
  \draw[->, very thick, WarnRed, dashed] (ff) -- (o);
  \node[font=\scriptsize\bfseries, text=WarnRed] at ($(ff)!0.5!(o)+(0,0.32)$) {blocked};
\end{tikzpicture}

  \vspace{0.7em}

  {\footnotesize X on the path $\Rightarrow$ fault effect \alert{cannot propagate} $\Rightarrow$ structurally \alert{untestable}}
  \note{Timing-critical paths cannot afford the scan-MUX delay, so the most timing-critical 10\% of flip-flops are kept non-scan. At one frame those FF outputs are unknown (X). When a fault's propagation path runs through such an FF output, the X blocks the effect from reaching an observable output, making the fault structurally untestable. Recovering these faults is the whole problem.}
\end{frame}
